\documentclass[9pt,compress,aspectratio=169]{beamer}

% ========================================================
% CUSTOMIZATIONS 
\newcommand{\authorname}{Your Name}
\newcommand{\shortauthor}{Y. Name}
\newcommand{\authoremail}{your.name@polymtl.ca}
\newcommand{\coauthors}{Martin Lévesque, Daniel Therriault, Frédérick Gosselin}
\newcommand{\institution}{Polytechnique Montréal}
\newcommand{\authoraffiliation}{Laboratory for Multiscale Mechanics}
\newcommand{\slidesettitle}{Awesome Presentation Title}
\newcommand{\shorttitle}{Presentation}
\newcommand{\slidesetsubtitle}{}
\newcommand{\logoimagepath}{logos/polytechnique_logo}
\newcommand{\logoimagepathlab}{logos/Logo-LM2-nowriting}
\newcommand{\lastimagepath}{logos/lm2_may2018.png}

% ==========================================================
%% PACKAGES
\usepackage[utf8]{inputenc}
%\usepackage[english]{babel}
\usepackage[english]{babel}
\usepackage{tikz}
\usepackage{adjustbox}
\usepackage[T1]{fontenc}
\usepackage{graphicx}
\usepackage{tikz}
\usepackage{textpos}
\usepackage{color}
\usepackage{etex}
\usepackage{rotating}
\usepackage{hyperref}
\usepackage{import}
\usepackage{media9}
\usepackage{subfig}
\DeclareCaptionLabelSeparator{horse}{:\quad} % change according to your needs
\usepackage{textcomp,amsfonts,stmaryrd,soulutf8}
\usepackage{xifthen}% provides \isempty test
\usepackage{longtable,ifthen,xspace,url}
\usepackage{verbatim,pdflscape}
\usepackage[babel=true]{csquotes} % style de guillements 
\usepackage[style=authortitle,backend=bibtex]{biblatex}
\usepackage{caption}
\captionsetup{font=scriptsize,labelfont=scriptsize}

\bibliography{Biblio} 
\graphicspath{{./images/}}

%% Maths
\usepackage{mathtools,mathabx,mathpazo}
\usepackage{amsmath,amssymb,amscd,mathrsfs,fixmath}
%fraction a/b
\newcommand{\rfrac}[2]{\raisebox{0.5ex}{$#1$}\slash\raisebox{-0.5ex}{$#2$}} 
% function f->...
\newcommand{\function}[5]{\begin{array}{crcl}
#1 \, : & #2 & \longrightarrow & #3 \\
    & #4 & \longmapsto & #5 \end{array}}
%\begin{equation}
%\function{d}{\mathbb{R}^3}{\mathbb{R}^3}{\boldsymbol\xi}{\boldsymbol{d}(\boldsymbol\xi) = \boldsymbol\eta}  
%\end{equation}

% ==========================================================
%BEGIN DOCUMENT
%\begin{document}
% ==========================================================
% DO NOT MODIFY THE NEXT LINE
% First slide + template
\usetikzlibrary{decorations.fractals}
\useoutertheme[subsection=false,shadow]{miniframes}
\useinnertheme{default}
\usefonttheme{serif}
\usefonttheme{professionalfonts}
\useoutertheme{infolines}

\setbeamersize{text margin left=0mm,text margin right=0mm} 

%POLYMTL COLORS
\definecolor{polyblue}{HTML}{41AAE6}
\definecolor{polygreen}{HTML}{8CC83C}
\definecolor{polyorange}{HTML}{FA961E}
\definecolor{polyred}{HTML}{B91E32}
\definecolor{polygray}{HTML}{A6A8AB}

%SIDEBAR
\usetheme[hideothersubsections]{Goettingen}
\setbeamersize{sidebar width right=0.12\paperwidth}
\makeatletter
\setbeamertemplate{sidebar canvas \beamer@sidebarside}[vertical shading][top=polyred!15,bottom=polyred!10]
\makeatother
\usepackage{etoolbox}
\patchcmd\insertverticalnavigation{\dohead}{\vskip-25pt\dohead}{}{}

%%% TABLE Packages
\usepackage{array,multirow,makecell}
\setcellgapes{1pt}
\makegapedcells
\usepackage{colortbl}
\newcolumntype{R}[1]{>{\raggedleft\arraybackslash }b{#1}}
\newcolumntype{L}[1]{>{\raggedright\arraybackslash }b{#1}}
\newcolumntype{C}[1]{>{\centering\arraybackslash }b{#1}}
%Epaisseur lignes horizontales tableaux
\newlength\epaisLigne 
 \newcommand\Ghline{\noalign{\global\epaisLigne\arrayrulewidth\global
                              \arrayrulewidth 3pt}
		    \hline\noalign{\global\arrayrulewidth\epaisLigne}} 
 \newcommand\Mhline{\noalign{\global\epaisLigne\arrayrulewidth\global
                              \arrayrulewidth 1.5pt}
		    \hline\noalign{\global\arrayrulewidth\epaisLigne}}

% BEAMER TEMPLATE
\setbeamertemplate{caption}[numbered]
\setbeamertemplate{navigation symbols}{}

% Table of Contents
\setbeamertemplate{section in toc}[circle]
\setbeamerfont{section in toc}{shape=\bfseries\scshape, size=\large}
\setbeamerfont{footnote}{size=\tiny}
\setbeamerfont{caption}{size=\scriptsize}
\setbeamercolor{section number projected}{fg=white,bg=polyred}
\setbeamerfont{section number projected}{series=\bfseries, size=\normalsize}
\setbeamertemplate{subsection in toc}[square]
\setbeamercolor{subsection number projected}{bg=polyorange}

\setbeamerfont{frametitle}{shape=\bfseries\scshape}
\setbeamerfont{frametitle}{size=\LARGE}
\setbeamerfont{framesubtitle}{shape=\normalfont\bfseries}
\setbeamerfont{framesubtitle}{size=\normalsize}
\setbeamerfont{block title}{shape=\normalfont\bfseries}
\setbeamerfont{block title}{size=\normalsize}

\setbeamercolor{lower separation line head}{bg=DeepSkyBlue4} 
\setbeamercolor{normal text}{fg=black} 
\setbeamercolor{alerted text}{fg=polyred} 
\setbeamercolor{example text}{fg=black} 
\setbeamercolor{frametitle}{fg=polyred} 
\setbeamercolor{framesubtitle}{fg=polyorange} 
\setbeamercolor{structure}{fg=black} 
\setbeamercolor{palette tertiary}{fg=black} 
\setbeamercolor{palette quaternary}{fg=black} 
\setbeamercolor{block title}{fg=polyblue}

% BIB RESOURCE
\addbibresource{Biblio.bib}

% HEADBARCOLOR FIRST SLIDE MASKED
\newcommand{\headbarcolor}{polyred}
\setbeamertemplate{headline}{% 
	\setbeamercolor{head1}{bg=\headbarcolor}
	\hbox{%
  \begin{beamercolorbox}[wd=.01\paperwidth,ht=0ex,dp=0ex,center]{head1}%
  \fontsize{5}{5}\selectfont
  \end{beamercolorbox}%
  }
  \vspace{-25ex}
}

% FOOT LINE / FOOTPAGE / FOOTNOTE
\newcommand{\footref}[1]{\textsuperscript{\ref{#1}}}
\let\Tiny=\tiny
\setbeamertemplate{footline}{
\begin{tiny}
\setbeamercolor{foot1}{fg=polyblue,bg=polyblue!10}
\setbeamercolor{foot2}{fg=polygreen,bg=polygreen!10}
\setbeamercolor{foot3}{fg=polyorange,bg=polyorange!10}
\setbeamercolor{foot4}{fg=polyred,bg=polyred!10}

% taken from theme infolines and adapted
  \leavevmode%
  \hbox{%
  \begin{beamercolorbox}[wd=.4\paperwidth,ht=2.25ex,dp=1ex,center]{foot1}%
	\shorttitle
  \end{beamercolorbox}%
  
  \begin{beamercolorbox}[wd=.2\paperwidth,ht=2.25ex,dp=1ex,center]{foot2}
	\shortauthor
  \end{beamercolorbox}%
  
  \begin{beamercolorbox}[wd=.25\paperwidth,ht=2.25ex,dp=1ex,center]{foot3}
	\authoremail
  \end{beamercolorbox}%
  
  \begin{beamercolorbox}[wd=.15\paperwidth,ht=2.25ex,dp=1ex,right]{foot4}%
	\insertframenumber{} / \inserttotalframenumber \hspace*{2ex} 
  \end{beamercolorbox}}%
  
  \vskip0pt%
\end{tiny}
\vskip2pt
}

% HIDE SIDEBAR FIRST SLIDE
\newenvironment{specialframe}
{
    \begingroup
    \advance\textwidth1.2cm % see beamerthemeGoettingen.sty for the number
    \hsize\textwidth
    \columnwidth\textwidth
    \begin{frame}[plain]
}
{
    \end{frame}
    \endgroup
}

% LAST SLIDE
\newenvironment{lastframe}
{
    \begingroup
    %\advance\textwidth1.6cm % see beamerthemeGoettingen.sty for the number
    \hsize\paperwidth
    \columnwidth\textwidth
    \begin{frame}[plain]
}
{
    \end{frame}
    \endgroup
}

% %%%%%%%%%%%%%%%%%%%%%%%%%%%%%%%
% PDF INFO

\hypersetup{
  pdftitle={\slidesettitle},
  pdfauthor={\authorname},
  pdfstartview={Fit},
  colorlinks=true,            % colorise les liens
  breaklinks=true,            % permet les retours à la ligne pour les liens trop longs
  urlcolor=black,             % couleur des hyperliens
  linkcolor=black,            % couleur des liens internes aux documents
  citecolor=black             % couleur des liens vers les references
}

% %%%%%%%%%%%%%%%%%%%%%%%%%%%%%%% 
% FIRST SLIDE TEMPLATE
%\title{\slidesettitle}
%\subtitle{\slidesetsubtitle}
%\author{\authorname}

\begin{document}
%\tracingall
%\nobibliography{Biblio}

{\setbeamertemplate{footline}{}

    \begin{specialframe}
        \begin{center}
            \includegraphics[width=1.05\textwidth]{./logos/poly_sky.png}
        \end{center}

        \vspace{-0.8cm}

        \begin{block}{}
            \begin{columns}
                \column{0.18\textwidth}
                    \begin{minipage}[t][.2\textheight][c]{\linewidth}
                    \begin{figure}
                        \flushleft
                        \includegraphics[height=12mm]{\logoimagepath}
                    \end{figure}
                    \end{minipage}    
                \column{0.1\textwidth}
                \column{0.68\textwidth}	
                    \begin{minipage}[c][.2\textheight][c]{\linewidth}
                        \centering
                        {\Large \textcolor{polyred}{\textbf{\scshape{\slidesettitle}}}}
                    \end{minipage}
                \end{columns}

                \vspace{-0.6cm}

                \begin{columns}
                    \column{0.22\textwidth}
                    \begin{figure}
                        \flushleft
                        \includegraphics[height=7mm]{\logoimagepathlab}
                    \end{figure}
                    \column{0.1\textwidth}
                    \column{0.68\textwidth}
                        \begin{columns}
                                \column{0.8\textwidth}
                                \flushleft
                                \normalsize{\textbf{\authorname}}, \coauthors \\
                                {\small \authoraffiliation, \institution}\\
                                {\footnotesize \authoremail}
                            \column{0.2\textwidth}
                                \begin{minipage}[b][.2\textheight][b]{\linewidth}
                                    \flushright
                                    {\footnotesize \textcolor{gray}{\insertdate}}
                                \end{minipage}
                        \end{columns}     
                \end{columns}
        \end{block}
    \end{specialframe}
}

% %%%%%%%%%%%%%%%%%%%%%%%%%%%%%%% 
% OTHER SLIDES TEMPLATE

\setbeamertemplate{headline}{% 
	\setbeamercolor{head1}{bg=\headbarcolor}
	 \hbox{%
  \begin{beamercolorbox}[wd=.01\paperwidth,ht=2.25ex,dp=11ex,center]{head1}%
  \fontsize{5}{5}\selectfont
  \end{beamercolorbox}%
  }
  \vspace{-15.5ex}
	\begin{flushright}
		\includegraphics[height=5mm]{\logoimagepath} \\
		\vspace{1mm}
		\includegraphics[height=5mm]{\logoimagepathlab}

	\end{flushright}
}
% greater line spacing for new slides
%\linespread{1} 

% %%%%%%%%%%%%%%%%%%%%%%%%%%%%%%% 

% DEFINE LAST SLIDE
\setbeamertemplate{endpage}{%
    \begin{lastframe}

        \begin{figure}
            \centering
            \includegraphics[height=50mm]{\lastimagepath}
        \end{figure}

    \vspace{-0.5cm}
    
    \begin{block}{Funding}

\vspace{3mm}

% Uncomment the lines for each logo that is required
\includegraphics[height=15mm]{\logoA} \hfill 
\includegraphics[height=18mm,page=1]{\logoB} \hfill 
\includegraphics[height=18mm,page=2]{\logoC} \hfill 
%\includegraphics[height=10mm]{\logoD} \hfill 
%\includegraphics[height=10mm]{\logoE} \hfill

    \end{block}

    \end{lastframe}
}

% IMPLEMENT LAST SLIDE
\AtEndDocument{\usebeamertemplate{endpage}}

% ======================= BEGIN SLIDES =====================
	
%Sets

\captionsetup{font=scriptsize, labelfont=scriptsize}


\section{Introduction}
%%%%%%%%%%%%%%%%%%%%%%%%%%%%%%%%%%%%%%%%%%%%%%%%%%%%%%%%%%%%%%%%%%%%%%
\begin{frame}

	\frametitle{Introduction}
	\framesubtitle{Fiber Reinforced Composites (FRCs)}

	\begin{center}
		\includegraphics[width=0.5\textwidth]{testimage.jpg} \\
		This is a test image with a footnote citation \footcite{CitekeyArticle} 
		
	\end{center}

\end{frame}
%%%%%%%%%%%%%%%%%%%%%%%%%%%%%%%%%%%%%%%%%%%%%%%%%%%%%%%%%%%%%%%%%%%%%%
\begin{frame}

	\frametitle{Test slide \#2}
	\framesubtitle{This slide contains equations and text}
	
	\hfill
	
	\begin{align*}
	f(x) &= x^2\\
	g(x) &= \frac{1}{x}\\
	F(x) &= \int^a_b \frac{1}{3}x^3
	\end{align*}
	
	\hfill
	
	Lorem ipsum dolor sit amet, consectetur adipiscing elit, sed do eiusmod tempor incididunt ut labore et dolore magna aliqua. Ut enim ad minim veniam, quis nostrud exercitation ullamco laboris nisi ut aliquip ex ea commodo consequat. Duis aute irure dolor in reprehenderit in voluptate velit esse cillum dolore eu fugiat nulla pariatur. Excepteur sint occaecat cupidatat non proident, sunt in culpa qui officia deserunt mollit anim id est laborum. \footnote{This is a footnote}
	
	\hfill
	
\end{frame}
%%%%%%%%%%%%%%%%%%%%%%%%%%%%%%%%%%%%%%%%%%%%%%%%%%%%%%%%%%%%%%%%%%%%%%
%=================================================================================================
%=================================================================================================



\end{document}
